\chapter{非负函数的积分}

\begin{solution}{simple-integral-refinement}
把 \(E_1,\ldots,E_m\) 细分为
\[
 A_\varepsilon=\bigcap_{\varepsilon_k=1}E_k
 \cap\bigcap_{\varepsilon_k=0}E_k^c,
 \qquad \varepsilon\in\{0,1\}^m.
\]
补集均相对于 \(X\) 取。删去空集后，这些集合两两不交并覆盖 \(X\)；
全零指标对应 \(X\setminus\bigcup_kE_k\)，其上的函数值为零。在
\(A_\varepsilon\) 上，\(\varphi\) 的值是
\(\sum_k a_k\varepsilon_k\)，故规范表示给出的积分为
\[
 \sum_\varepsilon\Bigl(\sum_ka_k\varepsilon_k\Bigr)
 \mu(A_\varepsilon)
 =\sum_ka_k\sum_{\varepsilon_k=1}\mu(A_\varepsilon)
 =\sum_ka_k\mu(E_k).
\]
所有项非负，有限和的换序在含 \(+\infty\) 时仍合法；
\(a_k=0,\mu(E_k)=\infty\) 时该项按约定为零。
\end{solution}

\begin{solution}{simple-integral-level-sets}
设规范表示为
\(\varphi=\sum_{j=1}^rb_j\mathbf1_{A_j}\)，不同正值按
\(0<b_1<\cdots<b_r\) 排列。函数
\(t\mapsto\mu(\{\varphi>t\})\) 是允许取无穷值的阶梯函数。
若没有正取值，则两端均为零。若某个 \(\mu(A_j)=+\infty\)，左端积分为无穷，
而 \(0\le t<b_j\) 时水平集测度也为无穷，故右端同为无穷。
以下设所有这些测度有限。逐点使用
\[
 b_j=\int_0^\infty\mathbf1_{\{b_j>t\}}\,dt
\]
并交换有限求和与积分，得到
\[
 \int_0^\infty\mu(\{\varphi>t\})\,dt
 =\sum_j\mu(A_j)b_j=\int\varphi\,d\mu.
\]
\end{solution}

\begin{solution}{nonnegative-integral-localization}
对简单函数 \(\varphi\le f\)，
\[
 \int_{E\cup F}\varphi=\int_E\varphi+\int_F\varphi.
\]
取一列简单函数 \(\varphi_n\uparrow f\)，则
\(\varphi_n\mathbf1_E\uparrow f\mathbf1_E\)，对 \(F,E\cup F\) 同理。
单调收敛定理和扩充实数中的加法连续性给出所求等式。
\end{solution}

\begin{solution}{nonnegative-integral-ae-invariance}
设 \(N=\{f\ne g\}\)，则 \(\mu(N)=0\)。任意非负简单函数
\(\varphi\le f\) 满足
\(\varphi\mathbf1_{N^c}\le g\)，并且
\(\int\varphi=\int\varphi\mathbf1_{N^c}\)，因为支撑在 \(N\) 上的简单
函数积分为零。因此 \(\int f\le\int g\)；交换 \(f,g\) 得等号。
\end{solution}

\begin{solution}{simple-approximation-alternative}
令
\[
 E_{n,k}=\{k2^{-n}\le f<(k+1)2^{-n}\},
 \quad 0\le k<2^{2n},
\]
并令 \(E_{n,*}=\{f\ge2^n\}\)。取
\[
 \varphi_n=\sum_{k=0}^{2^{2n}-1}k2^{-n}\mathbf1_{E_{n,k}}
 +2^n\mathbf1_{E_{n,*}}.
\]
各集合可测，故 \(\varphi_n\) 是简单函数。二进分割在下一层细分，向下
端点不减，所以 \(\varphi_n\le\varphi_{n+1}\le f\)。有限值点最终不被
截断且误差小于 \(2^{-n}\)；无穷值点上 \(\varphi_n=2^n\)。
\end{solution}

\begin{solution}{mct-indicator-continuity}
若 \(E_n\uparrow E\)，则
\(\mathbf1_{E_n}\uparrow\mathbf1_E\)。单调收敛定理给出
\[
 \mu(E)=\int\mathbf1_E
 =\lim_n\int\mathbf1_{E_n}=\lim_n\mu(E_n).
\]
单调收敛定理的证明只对简单函数使用了测度连续性；那一性质已在第
\ref{chap:measure-spaces} 章由可数可加性证明。因此这里是由已知测度性质
验证积分定理的特例，不是用积分
定理反过来证明其自身前提。
\end{solution}

\begin{solution}{mct-decreasing-counterexample}
在 \((\mathbb R,\mathcal L,m)\) 上取
\(f_n=\mathbf1_{[n,\infty)}\)。则 \(f_n(x)\downarrow0\) 对每个 \(x\)
成立，但 \(\int f_n\,dm=\infty\)。

相反，设非负可测函数 \(f_n\downarrow f\)，且
\(A=\int f_{n_0}\,d\mu<\infty\)。由 Markov 估计，对每个 \(k\ge1\)，
\(\mu(\{f_{n_0}=\infty\})\le A/k\)，所以这个可测集合为零集。
将 \(f\) 及所有 \(f_n\) 在其上置零，不改变积分。对 \(n\ge n_0\)，
此时 \(f_{n_0}\) 处处有限，故可定义
\(g_n=f_{n_0}-f_n\uparrow g=f_{n_0}-f\)。单调收敛给
\(\int g_n\uparrow\int g\)。非负加法公式给
\[
 A=\int g_n+\int f_n=\int g+\int f.
\]
各项均不超过有限数 \(A\)，于是只对有限实数作减法，得到
\(\int f_n=A-\int g_n\downarrow A-\int g=\int f\)。
若递减关系仅几乎处处成立，先在一个可测零集上把所有函数统一置零即可。
\end{solution}

\begin{solution}{nonnegative-series-tonelli}
在 \(\mathbb N\) 上取计数测度，令 \(f_m(n)=a_{mn}\)。非负函数逐项积分
公式给出
\[
 \sum_m\sum_na_{mn}
 =\int_{\mathbb N}\sum_m f_m(n)\,d\#(n)
 =\sum_n\sum_ma_{mn}.
\]
也可把两边都写成所有有限矩形部分和的上确界；非负性保证不同求和次序取到
同一个上确界。
\end{solution}

\begin{solution}{zero-integral-superlevel}
由 \(c\mathbf1_{\{f\ge c\}}\le f\)，
\[
 \int f\ge c\,\mu(\{f\ge c\})>0.
\]
在 \([0,1]\) 的 Lebesgue 测度上，取
\(f=\mathbf1_{\mathbb Q\cap[0,1]}\)。它在稠密集
\(\mathbb Q\cap[0,1]\) 上为正，但非零点集零测，故积分为零。这里非零点集
的闭包却是整个区间，不能把它与拓扑意义的支集混同。稠密性是拓扑
性质，并不保证函数在正测度集合上具有正下界。
\end{solution}

\begin{solution}{density-measure}
\(\nu(\varnothing)=0\)。若 \(E_n\) 两两不交，则
\[
 f\mathbf1_{\cup_nE_n}=\sum_nf\mathbf1_{E_n}.
\]
非负函数的可数逐项积分公式给出
\[
 \nu\Bigl(\bigcup_nE_n\Bigr)=\sum_n\nu(E_n),
\]
故 \(\nu\) 是测度。若 \(\mu(E)=0\)，则
\(f\mathbf1_E=0\) 几乎处处，所以 \(\nu(E)=0\)，即 \(\nu\ll\mu\)。
\end{solution}

\begin{solution}{cavalieri-simple-to-general}
简单函数情形由习题 \ref{ex:simple-integral-level-sets}。取
\(\varphi_n\uparrow f\)。对每个 \(t\)，集合
\(\{\varphi_n>t\}\uparrow\{f>t\}\)，故测度连续性给出
\[
 \mu(\{\varphi_n>t\})\uparrow\mu(\{f>t\}).
\]
分别在 \(X\) 和 \([0,\infty)\) 上使用单调收敛定理：
\[
 \int_Xf\,d\mu
 =\lim_n\int_X\varphi_n\,d\mu
 =\lim_n\int_0^\infty\mu(\{\varphi_n>t\})\,dt
 =\int_0^\infty\mu(\{f>t\})\,dt.
\]
\end{solution}

\begin{solution}{integral-from-functional-axioms}
取定理 \ref{thm:nonnegative-simple-approximation} 中的
\(\varphi_n\uparrow f\)。由假设，
\[
 I(f)=\lim_nI(\varphi_n)=\lim_n\int\varphi_n\,d\mu=\int f\,d\mu.
\]
最后一个等号是单调收敛定理。事实上只要 \(I\) 在递增简单函数列上具有
从下连续性，以上论证已经成立。
\end{solution}
